\section{Resultados}\label{sec:Resultados}

A metodologia foi avaliada em 23 processadores RISC-V \textit{open-source}
para os quais os resultados da versão atual do Processor CI Inspector foram
confirmados por revisão independente do RTL, da configuração selecionada e,
quando disponível, da documentação do projeto. O Kronos foi mantido como um
estudo de caso adicional de divergência, pois sua classificação como
\textit{pipeline} está correta, embora a profundidade temporal estimada não
coincida com os três estágios declarados pelo projeto.

A ferramenta foi empregada para detectar automaticamente os seguintes
parâmetros:

\begin{description}
  \item[Tipo de licença:] identificação das licenças encontradas nos arquivos
  do projeto;
  \item[Tamanho da palavra:] determinação da largura arquitetural a partir do
  banco de registradores;
  \item[Modelo de execução:] classificação do \textit{datapath} como
  monociclo, multiciclo ou \textit{pipeline};
  \item[Profundidade estimada:] estimativa temporal calculada a partir da
  latência modal entre a busca e a conclusão arquitetural da instrução;
  \item[Latência de instrução:] número modal de ciclos entre a busca da
  instrução e a observação de seu \textit{commit};
  \item[Linguagem de descrição de hardware:] identificação da linguagem HDL
  predominante na implementação selecionada.
\end{description}

Entre os 23 resultados confirmados, 17 processadores foram implementados em
Verilog, quatro em SystemVerilog e dois em VHDL. Quanto ao modelo de execução,
sete foram classificados como \textit{pipeline}, 13 como multiciclo e três
como monociclo. Vinte e um implementam arquiteturas RV32 e dois implementam
RV64. A Tabela \ref{tab:resultados} apresenta esses resultados e inclui o
Kronos, identificado explicitamente como caso de divergência.

\begin{table*}[t]
\centering
\begin{threeparttable}
\resizebox{\textwidth}{!}{%
\begin{tabular}{|l|l|c|c|c|c|c|c|}
  \hline
  \textbf{Processador} & \textbf{Licença detectada} & \textbf{XLEN} &
  \textbf{Tipo detectado} & \textbf{$D_{\mathrm{est}}$} &
  \textbf{$L_{\mathrm{mode}}$} & \textbf{$D_{\mathrm{ref}}$} &
  \textbf{HDL} \\
  \hline
  Baby-Risco-5 & CERN-OHL-P-2.0, MIT e domínio público & 32 & Multiciclo & -- & 6 & -- & Verilog \\
  \hline
  RPU & Apache 2.0 & 32 & Multiciclo & -- & 4 & -- & VHDL \\
  \hline
  Risco-5 & CERN-OHL-P-2.0, MIT e domínio público & 32 & Multiciclo & -- & 8 & -- & Verilog \\
  \hline
  arRISCado & GPLv3 & 32 & Pipeline & 4 & 3 & N/D\tnote{1} & Verilog \\
  \hline
  clarvi & Não detectada & 32 & Pipeline & 6 & 5 & 6 & SystemVerilog \\
  \hline
  Uriscv & Apache 2.0 & 32 & Multiciclo & -- & 3\tnote{2} & -- & Verilog \\
  \hline
  cv32e40x & Solderpad & 32 & Pipeline & 4 & 3 & 4 & SystemVerilog \\
  \hline
  cve2 & Apache 2.0 & 32 & Pipeline & 2 & 1 & 2 & SystemVerilog \\
  \hline
  Fedar F1 & MIT & 64 & Pipeline & 5 & 4 & 5 & Verilog \\
  \hline
  harv & MIT & 32 & Multiciclo & -- & 3 & -- & VHDL \\
  \hline
  Mriscv & Licença customizada e MIT & 32 & Multiciclo & -- & 4 & -- & Verilog \\
  \hline
  Nerv & Não detectada & 32 & Monociclo & -- & 2 & -- & SystemVerilog \\
  \hline
  Picorv32 & Licença customizada & 32 & Multiciclo & -- & 5 & -- & Verilog \\
  \hline
  riscado-v & BSD 3-Clause & 32 & Multiciclo & -- & 5 & -- & Verilog \\
  \hline
  riscv & BSD 3-Clause & 32 & Pipeline & 6 & 5 & N/D\tnote{1} & Verilog \\
  \hline
  riscv\_pipeline & MIT & 32 & Pipeline & 4 & 3 & N/D\tnote{1} & Verilog \\
  \hline
  Riskow & Apache 2.0 & 32 & Multiciclo & -- & 5 & -- & Verilog \\
  \hline
  RVX & MIT & 32 & Monociclo & -- & 2 & -- & Verilog \\
  \hline
  SERV & Licença customizada & 32 & Multiciclo & -- & 14\tnote{3} & -- & Verilog \\
  \hline
  simple\_riscv\_cpu & Licença customizada e LGPLv3 & 32 & Multiciclo & -- & 6 & -- & Verilog \\
  \hline
  RISCV\_Single\_Cycle & Não detectada & 32 & Monociclo & -- & 0 & -- & Verilog \\
  \hline
  mystic\_riscv64 & MIT & 64 & Multiciclo & -- & 22\tnote{4} & -- & Verilog \\
  \hline
  riscv-core & MIT & 32 & Multiciclo & -- & 11\tnote{4} & -- & Verilog \\
  \hline
  Kronos\tnote{5} & Apache 2.0 & 32 & Pipeline & 6 & 5 & 3 & SystemVerilog \\
  \hline
\end{tabular}%
}

\begin{tablenotes}
\footnotesize
\item[1] A documentação ou o RTL confirma a presença de \textit{pipeline}, mas
não declara uma profundidade única que possa ser usada como referência.
\item[2] No Uriscv, quatro das seis amostras apresentaram três ciclos e duas
apresentaram 36 ciclos; a tabela mostra a moda, com intervalo de 3--36 ciclos.
\item[3] No SERV, a latência modal foi de 14 ciclos, observada em três das seis
amostras; o intervalo completo foi de 14--34 ciclos.
\item[4] Foram observadas cinco amostras válidas nesses processadores.
\item[5] O Kronos é apresentado como estudo de divergência e não integra o
conjunto dos 23 resultados confirmados.
\end{tablenotes}

\caption{Resultados da versão atual do Processor CI Inspector para os 23 casos
confirmados e para o Kronos, mantido como estudo de divergência.
$D_{\mathrm{est}}$ representa a estimativa temporal do Inspector,
$L_{\mathrm{mode}}$ a latência modal entre busca e \textit{commit}, e
$D_{\mathrm{ref}}$ a profundidade arquitetural verificada no RTL ou na
documentação.}
\label{tab:resultados}
\end{threeparttable}
\end{table*}

O termo ``resultado confirmado'' indica concordância entre a saída do
Inspector e a referência manual para XLEN, linguagem HDL e modelo de execução.
A licença detectada e a latência observada são apresentadas como resultados da
ferramenta, mas não foram utilizadas como critérios de inclusão nesse conjunto.
No caso de processadores com profundidade explicitamente documentada, a
comparação também considera esse valor.

Entre os sete processadores com \textit{pipeline} no conjunto confirmado,
quatro possuem profundidade fixa explicitamente documentada: clarvi, cv32e40x,
cve2 e Fedar F1. As estimativas obtidas para esses quatro processadores foram,
respectivamente, 6, 4, 2 e 5 estágios, coincidindo com as respectivas
referências. Para arRISCado, riscv e riscv\_pipeline, a presença de
\textit{pipeline} foi confirmada, mas não se adotou uma profundidade estrutural
única como referência; portanto, os valores 4, 6 e 4 devem ser interpretados
somente como estimativas temporais.

Os resultados também confirmaram 13 implementações multiciclo e três
implementações monociclo. A classificação do RVX como monociclo merece
destaque, pois seu RTL executa a decodificação, a operação combinacional e a
seleção do resultado entre atualizações do estado arquitetural. A antiga
classificação do RVX como multiciclo não foi mantida. As latências de dois
ciclos observadas no RVX e no Nerv não alteram essa classificação, pois podem
incluir a entrega registrada da instrução e o ponto em que o estado
arquitetural é observado.

\subsubsection{Kronos como Estudo de Divergência}

O Kronos é documentado como um processador RV32 de três estágios em ordem. O
Inspector identificou corretamente a execução em \textit{pipeline}, com uma
conclusão arquitetural por ciclo, mas observou latência modal de cinco ciclos.
Pela convenção atual,

\begin{equation}
    D_{\mathrm{est}} = L_{\mathrm{mode}} + 1,
\end{equation}

o que produziu uma estimativa de seis estágios. A descoberta de sinais internos
encontrou uma cadeia coerente contendo os papéis de \textit{frontend} e
\textit{execute}, além da interface de escrita do banco de registradores. No
entanto, o estágio de memória não foi marcado como observável e a busca de
candidatos foi truncada pelos limites da análise. Assim, o estado
\texttt{confirmed} da interface indica que há evidência suficiente para
análises funcionais subsequentes, mas não que todos os três estágios físicos
tenham sido enumerados.

A divergência do Kronos não indica erro na documentação do processador. Ela
mostra que a profundidade atualmente apresentada pelo Inspector é uma
estimativa temporal entre os pontos observados de busca e conclusão, podendo
incorporar ciclos da interface de instruções, do \textit{wrapper} ou do ponto
de observação do \textit{writeback}. Portanto, essa estimativa deve ser
distinguida da profundidade arquitetural declarada pelo projeto.

A detecção de licenças também deve ser interpretada como um levantamento dos
avisos encontrados no repositório, e não necessariamente como a licença única
do núcleo. Repositórios como Baby-Risco-5 e Risco-5 contêm arquivos associados
a mais de uma licença. Em outros casos, o texto não foi associado
automaticamente a um identificador conhecido: o Nerv e o
RISCV\_Single\_Cycle foram apresentados como ``Não detectada'', enquanto o
Picorv32, documentado pelo projeto como ISC, foi classificado genericamente
como licença customizada. Esses casos evidenciam limitações do reconhecedor de
licenças e não invalidam a caracterização do \textit{datapath}.

A ferramenta foi integrada ao ambiente de integração contínua do Processor CI,
permitindo que a inspeção seja executada quando um processador é adicionado ou
atualizado. Os resultados são armazenados em JSON para permitir tanto revisão
manual quanto consumo por outras ferramentas. A Figura \ref{fig:json} apresenta
um trecho da saída atual para o cve2, um dos casos cuja profundidade estimada
coincidiu com a documentação.

\begin{figure}[t]
    \centering
\begin{lstlisting}[
    language=json,
    firstnumber=1,
    captionpos=b
]
{
  "cve2": {
    "license_types": ["Apache 2.0"],
    "bits": 32,
    "language": "SystemVerilog",
    "multicycle": false,
    "pipeline": {
      "depth_estimate": 2,
      "depth_estimate_source":
        "write_interface_timing_corrected",
      "raw_depth_estimate": 3
    },
    "cycle": {
      "fetch_to_commit_latencies": [1, 1, 1, 1, 1, 1],
      "commit_intervals": [1, 1, 1, 1, 1]
    }
  }
}
\end{lstlisting}
\caption{Trecho da saída JSON produzida para o cve2.}
\label{fig:json}
\end{figure}
